\documentclass[12pt,letterpaper]{article}
\usepackage[utf8]{inputenc}
\usepackage[spanish]{babel}
\usepackage{amsmath}
\usepackage{amsfonts}
\usepackage{amssymb}
\usepackage[margin=2.5cm]{geometry}

\begin{document}

\begin{center}
    \Large \textbf{Evaluación de Examen 2: Unidad II} \\
    \large Física para Ciencias de la Salud (005-1713) \\
    \vspace{0.5cm}
    \textbf{Estudiante:} Krischel Bravo \quad \textbf{C.I.:} 32.449.538
    \vspace{0.3cm} \\
    \large \textbf{Calificación Final: 10/10 puntos}
\end{center}

\vspace{0.5cm}

\section*{Análisis de Resultados}

\subsection*{1. Cinemática (Aceleración media)}
\textbf{Procedimiento y resultado:} Extrae los datos correctamente y plantea la ecuación de aceleración media ($a = \frac{v_f - v_0}{t}$). Realiza la sustitución adecuada y el cálculo aritmético arroja el resultado analítico exacto de $4 \text{ m/s}^2$ con sus unidades correctas en el Sistema Internacional. \\
\textbf{Veredicto:} \textbf{Correcto} (2/2 pts).

\subsection*{2. Dinámica (Leyes de Newton)}
\textbf{Procedimiento y resultado:} Muestra un despeje directo a partir de la Segunda Ley de Newton ($a = \frac{F}{m}$). Sustituye los valores dividiendo $150 \text{ N}$ entre $25 \text{ kg}$, logrando un resultado perfecto de $6 \text{ m/s}^2$. \\
\textbf{Veredicto:} \textbf{Correcto} (2/2 pts).

\subsection*{3. Trabajo Mecánico}
\textbf{Procedimiento y resultado:} Plantea la fórmula general del trabajo ($W = F \cdot d \cdot \cos(\theta)$). Sustituye los valores numéricos correspondientes y efectúa la multiplicación ($400 \text{ N} \cdot 0,8 \text{ m} \cdot 1$), obteniendo de forma exacta $320 \text{ J}$. \\
\textbf{Veredicto:} \textbf{Correcto} (2/2 pts).

\subsection*{4. Energía Potencial}
\textbf{Procedimiento y resultado:} Extrae de manera estructurada los datos del enunciado (masa, altura y gravedad) y aplica a la perfección la fórmula de la energía potencial gravitatoria ($E_p = m \cdot g \cdot h$). El producto matemáticamente exacto es $17,64 \text{ Joules}$. \\
\textbf{Veredicto:} \textbf{Correcto} (2/2 pts).

\subsection*{5. Potencia Mecánica}
\textbf{Procedimiento y resultado:} Extrae correctamente los datos y plantea la fórmula de potencia mecánica ($P = \frac{W}{t}$). Ejecuta la sustitución de $75 \text{ J} / 60 \text{ s}$ logrando el resultado perfecto de $1,25 \text{ Watts}$. \\
\textbf{Veredicto:} \textbf{Correcto} (2/2 pts).

\vspace{0.5cm}
\section*{Comentarios Generales y Sugerencias}
¡Felicidades por obtener la nota máxima en la evaluación de la Unidad II! El examen evidencia un excelente dominio de los conceptos físicos y de cálculo matemático.
\begin{itemize}
    \item Se valora muy positivamente la estructura de resolución de la estudiante, mostrando en diferentes zonas de la hoja el desarrollo y el planteamiento de la fórmula original utilizada; esto facilita enormemente el seguimiento analítico en ciencias.
    \item Existe un buen uso del análisis dimensional, indicando correctamente cada unidad del Sistema Internacional junto al resultado (Joules, Watts, $\text{m/s}^2$).
    \item Como único y pequeño detalle de ortografía técnica, se recomienda escribir ``Watts'' o emplear su símbolo oficial ``W'' en mayúscula, en lugar de ``w'' minúscula, ya que proviene de un nombre propio.
    \item ¡Sigue con este fantástico nivel de dedicación!
\end{itemize}

\end{document}